\section{The benchmark}
\label{sec:benchmark}

\subsection{The task, the contract, and the scoring}
\label{sec:task}

The agent receives a frozen, versioned problem statement
(Appendix~\ref{app:tasktext}) giving: the forward problem---fixed-boundary
equilibria of D-shaped plasmas in five parameters ($\rzero$, $a$, $\kappa$,
$\delta$, $\Ip$) over stated ranges at the solver's default profiles; a campaign
structure with a hard budget---at most 150 initial space-filling solves then up
to three adaptive rounds of exactly 100 new solves each, failures included, with
each round's points required to be chosen from the current model or data, so a
one-shot design is explicitly insufficient; a stopping criterion---stop early
once validation error reaches $\le 0.30\times$ the mean-predictor baseline, on a
validation set kept separate from the test set; a held-out test set of at least
20 solves at a recorded seed, never used for training, selection or acquisition;
and a requirement to report whether the adaptive sampling actually helped, ``or
an honest statement that it did not''. Four process gates are mandatory and
identical for every condition: a \emph{meshing milestone}, a \emph{pilot gate}
($\ge 20$ solves at $\ge 50\%$ success before the rounds begin), a
\emph{storage-validation gate} (a solve counts only after its stored artifact is
re-read from disk and confirmed finite), and a \emph{deliverable self-test} (run
every documented command in a fresh process before declaring done).

\paragraph{The deliverable contract.}
Every condition emits the same artifacts, checked mechanically. The nine checks
(Appendix~\ref{app:gates}) require the three artifacts to exist and be
non-empty, \code{report.json} to parse and carry the required keys,
\code{predict.py} to run to completion as a subprocess on sample parameters,
and its output to have the right shape on the right grid. One
further gate applies at \Lthree{} only---a static audit that no workspace file
imports the platform library---so \Lthree{} runs face nine gates and \Ltwo{}
eight. \code{contract.passed} is the flat \textsc{and} of them: no weighting, no
partial credit. The scoring subprocess is stripped of provider credentials, so
\code{predict.py} must work offline, and runs in its own process group so a
timeout reaps the solver's children. Critically, \code{report\_keys}
\textbf{checks key \emph{presence}, never correctness}: self-reported metrics
are recorded verbatim and never cross-checked, which is what makes the honesty
measurement of Section~\ref{sec:honesty} possible.

\paragraph{Independent scoring.}
Predictions are scored against a frozen test set of 60 parameter vectors (seed
\code{20260809}, all 60 converged) solved once and stamped with provenance---
solver version, the SHA-256 of both the parameter file and the grid, and a
scoring epoch---so a score is only ever pooled with scores computed the same
way. The archive is never regenerated. The metric is per-sample relative $L_2$
over the finite ground-truth points, and \textbf{a NaN prediction at a finite
ground-truth point is scored as a full-magnitude miss}, never skipped, since
otherwise an all-NaN prediction would score as perfect
(Appendix~\ref{app:scoring}). Accuracy is quoted against the trivial
mean-flux-map predictor on the same set, which scores $0.4933$.

\subsection{The access-level ladder, and the substrates}
\label{sec:ladder}

Four levels, all emitting the identical contract and scored on the identical
frozen set: \Lzero{} \emph{scripted}, the pre-written pipeline with seeded
heuristic decisions and zero LLM calls; \Lone{} \emph{structured}, the same
pipeline with an LLM making only typed decisions; \Ltwo{}
\emph{library-assisted}, where the agent writes the glue code and may import the
platform library; and \Lthree{} \emph{from-scratch}, where it writes everything
and the import is forbidden and audited. \Lzero{} is the golden baseline---same
seed, same decisions, no model---and having a zero-LLM condition scored on the
same frozen set is what lets us say how much of a score is the task and how much
is the agent. Four agent substrates implement one adapter each behind the same
interface: \hn{ursa} (a plan/execute agent pair on LangGraph), \hn{dspy} (plan
$\rightarrow$ per-step ReAct $\rightarrow$ review $\rightarrow$ fix), and
\hn{pi} and \hn{cursor} (CLI coding agents); adding a substrate is one adapter
module and one registry entry.

\paragraph{Is the task hard for the right reasons?}
\label{sec:characterisation}
Three facts from a direct sweep independent of any agent, detailed in
Appendix~\ref{app:characterisation}. The \emph{best model family changes with
$N$}---polynomial ridge up to $N = 200$, kernel ridge from $400$---so model
selection depends on the data the agent chose, and the 450-solve bound sits
where that still matters. At matched budget space-filling beats uniform-random
by $21\%$ at $N = 100$ and $6\%$ by $1600$, while clustered sampling is
catastrophic everywhere. And the surrogate \emph{collapses out of
distribution}---trained on half the parameter space and tested on the other,
ratio $0.94$ against $0.38$--$0.40$ within-distribution---so an agent that
samples badly does not lose a few percent; it produces a model that fails
completely off-distribution and, evaluating itself on its \emph{own} test set,
may never notice.
